% Faithful mathematical transcription of official CCC2022 article21. Original author numbering retained.
\section*{Necessary original definitions and main result}
Let $\operatorname{sign}:\mathbb R\to\mathbb R$ be defined as $\operatorname{sign}(z)=1$ if $z\geq0$ and $0$ otherwise.
\paragraph{Definition 1.} For an integer $d>0$, a degree $d$ PTF $f:\mathbb R^n\to\{0,1\}$ is a function of the form $f(x)=\operatorname{sign}(p(x))$, where $p:\mathbb R^n\to\mathbb R$ is a polynomial of degree at most $d$.
Our goal is to design a PRG that takes few random bits and outputs a high-dimensional vector whose distribution is indistinguishable from a standard multivariate gaussian by any low-degree PTF. Specifically:
\paragraph{Definition 2.} A function $G:\{0,1\}^r\to\mathbb R^n$ is a pseudorandom generator for degree $d$ PTFs with error $\varepsilon$ if for every degree at most $d$ PTF $f:\mathbb R^n\to\{0,1\}$,
\[\left|\mathbb P_{y\in_u\{0,1\}^r}(f(G(y))=1)-\mathbb P_{x\sim N(0,1)^n}(f(x)=1)\right|\leq\varepsilon.\]
We call $r$ the seed length of the generator and say $G$ $\varepsilon$-fools degree $d$ PTFs with respect to the gaussian distribution.\footnote{Here, and henceforth, $y\in_u S$ denotes a uniformly random element from a multi-set $S$, and $N(0,1)$ represents the standard univariate gaussian distribution of variance $1$.} We say $G$ is explicit if its output can be computed in time polynomial in $n$.
\paragraph{Theorem 3 (PRG for PTFs).} There exist constants $c,C$ such that for all $\varepsilon>0$ and $d\geq1$, there exists an explicit PRG that $\varepsilon$-fools $n$-variate degree $d$ PTFs with respect to the gaussian distribution with seed length $r(n,d,\varepsilon)=Cd^c\log(n/\varepsilon)\log(1/\varepsilon)/\varepsilon^2$.
\paragraph{Definition 4.} Given a function $p:\mathbb R^n\to\mathbb R$ and $x\in\mathbb R^n$, and a restriction parameter $\lambda\in(0,1)$, let $p_{x,\lambda}:\mathbb R^n\to\mathbb R$ be the function $p_{x,\lambda}(y)=p(\sqrt{1-\lambda}x+\sqrt\lambda y)$.\footnote{As the value of $\lambda$ will often be clear, we will often in fact just use $p_x$ for brevity.}
Intuitively, we can view $p_{x,\lambda}$ as a restriction where $(1-\lambda)$-fraction of the variance is already fixed. (Note that for independent $x,y\sim N(0,1)^n$, $\sqrt{1-\lambda}x+\sqrt\lambda y$ is distributed as $N(0,1)^n$.)
For a function $p:\mathbb R^n\to\mathbb R$, let $\|\nabla^kp(x)\|^2$ denote the sum of squares of all partial derivatives of $p$ of order $k$ at $x$. That is, $\|\nabla^kp(x)\|$ is the Frobenius norm of the tensor of $k$'th order partial derivatives of $p$. We show that for any degree $d$ polynomial $p$, the Frobenius-norm of the $k$'th order derivatives are comparable to the $(k-1)$'th order derivatives on a random gaussian input with high probability:
\paragraph{Lemma 6.} For any degree-$d$ polynomial $p:\mathbb R^d\to\mathbb R$, and $x\sim N(0,1)^n$, the following holds with probability at least $1-\delta$:
\[\|\nabla^kp(x)\|\leq O(d^3/\delta)\|\nabla^{k-1}p(x)\|,\text{ for all }1\leq k\leq d.\tag{1}\]
\paragraph{Definition 10.} For a polynomial $p:\mathbb R^n\to\mathbb R$ of the form $p(y):=\sum_\alpha\widehat p(\alpha)h_\alpha(y)$, define its hypervariance, $\operatorname{HyperVar}_R(\ )$, and normalized hypervariance, $H_R(\ )$, as
\[\operatorname{HyperVar}_R(p):=\sum_{\alpha\neq0}\widehat p(\alpha)^2R^{2|\alpha|},\qquad H_R(p):=\frac{\operatorname{HyperVar}_R(p)}{\widehat p(0)^2}.\]
Intuitively, if the normalized hypervariance $H_R(p)$ of a polynomial is small for a large $R$, then it means that the weights of the higher-order Hermite coefficients of $p$ have a geometric decay.
\section*{3 Preliminaries}
\paragraph{The pseudorandom generator construction: idealization vs. discretization.} Following [7] and [14], we analyze the idealized pseudorandom distribution
\[Z=\frac1{\sqrt L}\sum_{i=1}^LY_i,\]
where each $Y_i\in\mathbb R^n$ is a $k$-moment-matching gaussian (that is, $\mathbb E[p(Y_i)]=\mathbb E_{x\sim N(0,1)^n}[p(x)]$ for all polynomials $p:\mathbb R^n\to\mathbb R$ of degree at most $k$).
Suppose that, for any such $Z$ with parameters $(L,k)$, it is the case that $Z$ fools degree-$d$ PTFs with error $\varepsilon=\varepsilon(L,k,d)$. Then, it is shown in [7] how to obtain a small-seed length PRG (in the sense of Definition 2) by providing a specific instantiation and discretization of this construction.
\paragraph{Theorem 12 ([7], implicit in Section 6).} Suppose a $Z$ as above with parameters $(L,k)$ fools degree $d$-PTFs with error $\varepsilon=\varepsilon(L,k,d)$. Then, there is an explicit, efficiently computable PRG with seed length $O(dkL\log(ndL/\varepsilon))$ that $(2\varepsilon)$-fools degree $d$ PTFs.
\paragraph{Hermite polynomials.} To argue about polynomials which are not necessarily multilinear, we need some simple facts concerning Hermite polynomials. For our purposes, Hermite polynomials are simply a convenient choice of polynomial basis which have nice properties (in particular being orthonormal) with respect to gaussian inputs. For a more detailed background on Hermite polynomials and their use for analyzing functions over gaussian space, see [13, Ch. 11].
One concrete way to define the Hermite polynomials is the following:
\begin{itemize}
\item For the univariate polynomials, the degree-$m$ ``Probabilist's'' Hermite polynomial is the $m$-th coefficient of the generating function
\[e^{st-\frac12s^2}=\sum_{m\geq0}H_m(t)s^m.\]
\item We define the degree-$m$ univariate Hermite polynomial by the normalization
\[h_m(t):=\frac1{\sqrt{m!}}H_m(t).\]
\item For a multi-index $\alpha\in\mathbb N^n$, we define the multivariate Hermite polynomial $h_\alpha:\mathbb R^n\to\mathbb R$ via the product
\[h_\alpha(x):=\prod_{i=1}^nh_{a_i}(x_i).\]
\end{itemize}
We record some basic properties of this particular choice of polynomial basis. The final two properties say that the Hermite basis is orthonormal with respect to correlation under the standard gaussian distribution -- this is the reason for our choice of normalization.
\begin{itemize}
\item The set $\{h_\alpha(x):|\alpha|\leq d\}$ is a basis for real polynomials in $n$ variables of degree $\leq d$.
\item $h_0$ is the constant polynomial $h_0\equiv1$.
\item For multi-indicies $\alpha\in\{0,1\}^n$, $h_\alpha(x)$ is simply the monomial $\prod_{i:a_i=1}x_i$.
\item For $x\sim N(0,1)^n$, and distinct multi-indices $\alpha\neq\beta$, $\mathbb E_xh_\alpha(x)h_\beta(x)=0$.
\item For $x\sim N(0,1)^n$, and any multi-index $\alpha$, $\mathbb E_xh_\alpha(x)^2=1$.
\end{itemize}
\paragraph{Gaussian noise operator.} We recall the definition of the noise operator $U_\rho$, which here we regard as an operator on real polynomials in $n$ variables (see [13, Ch. 11] for background and a more general viewpoint). For a polynomial $f:\mathbb R^n\to\mathbb R$ and a parameter $\rho\in[0,1]$, the action of $U_\rho$ on $f$ is specified by
\[(U_\rho f)(x):=\mathbb E_{Z\sim N(0,1)^n}f(\rho x+\sqrt{1-\rho^2}Z).\]
An important feature of the Hermite basis is that the noise operator acts on it diagonally (see [13, Ch. 11]): $U_\rho h_\alpha(x)=\rho^{|\alpha|}h_\alpha(x)$. Thus, if $f$ is a degree-$d$ polynomial given in the Hermite basis as
\[f(x)=\sum_{|\alpha|\leq d}\widehat f(\alpha)h_\alpha(x),\]
then we can express the result of the noise operator applied to $f$ explicitly as
\[U_\rho f(x)=\sum_{|\alpha|\leq d}\widehat f(\alpha)\rho^{|\alpha|}h_\alpha(x).\]
\paragraph{Higher moments and hypercontractivity.} Fix a polynomial $f(x):=\sum_{|\alpha|\leq d}\widehat f(\alpha)h_\alpha(x)$. For an even natural number $q\geq2$, we write the gaussian $q$-norm of $f$ as
\[\|f\|_q:=\left(\mathbb E_{x\sim N(0,1)^n}f(x)^q\right)^{1/q}.\]
We wish to be able to bound this quantity in terms of the magnitudes of the Hermite coefficients of $f$, $\widehat f(\alpha)$. For this purpose, we extend the definition of $U_\rho$ also to $\rho>1$ by its action on the Hermite basis: $U_\rho h_\alpha(x)=\rho^{|\alpha|}h_\alpha(x)$. With this notation, we can express the well-known $(q,2)$-hypercontractive inequality [13, Ch. 9,11] as
\[\|f\|_q\leq\|U_{\sqrt{q-1}}f\|_2,\]
which is quite convenient for us, as we can use orthonormality of the Hermite basis to explicitly compute
\[\|U_{\sqrt{q-1}}f\|_2^2=\sum_{|\alpha|\leq d}(q-1)^{|\alpha|}\widehat f(\alpha)^2\leq\sum_{|\alpha|\leq d}q^{|\alpha|}\widehat f(\alpha)^2.\]
\paragraph{Gaussian restrictions and derivatives on the Hermite basis.} Besides the effect of the noise operator, it will also be important to understand the effect of two further operations on polynomials:
\begin{itemize}
\item The derivative map, $f(y)\mapsto\partial^\alpha f(y)$.
\item The gaussian restriction at $x$, $f(y)\mapsto f(\sqrt{1-\lambda}x+\sqrt\lambda y)$.
\end{itemize}
In particular, we are concerned with how these operations affect the Hermite coefficients of a polynomial; ultimately, our goal will be to develop a ``Hermite-basis analogue'' of the Taylor expansion which can be applied to expand $f(\sqrt{1-\lambda}x+\sqrt\lambda y)$ as a function of $y$. We start by computing the effect of these two operations on univariate Hermite polynomials, and then on the full multivariate Hermite basis, and finally on a general polynomial $f(x)$ expressed in the Hermite basis.
\paragraph{Proposition 14.} For univariate Hermite polynomials, we have the identities
\begin{itemize}
\item $\frac{\partial^k}{\partial t^k}h_m(t)=\sqrt{\frac{m!}{(m-k)!}}h_{m-k}(t)$,
\item $h_m(\sqrt{1-\lambda}x+\sqrt\lambda y)=\sum_{k=0}^m\sqrt{\binom mk}(1-\lambda)^{(m-k)/2}\lambda^{k/2}h_{m-k}(x)h_k(y)$.
\end{itemize}
\paragraph{Proof.} The first of these identities is standard (see e.g. [13, Ex. 11.10]); we provide a proof of the second.
The second identity can be proved by considering the generating function
\[e^{st-\frac12s^2}=\sum_m\sqrt{m!}h_m(t)s^m,\]
and comparing the coefficient of $s^m$ on both sides of
\[e^{s(\sqrt{1-\lambda}x+\sqrt\lambda y)-\frac12s^2}=e^{(s\sqrt{1-\lambda})x-\frac12(s\sqrt{1-\lambda})^2}\cdot e^{(s\sqrt\lambda)y-\frac12(s\sqrt\lambda)^2}.\]\hfill$\triangleleft$
The corresponding identities for multivariate Hermite polynomials follow easily from above.
\paragraph{Proposition 15.} We have
\begin{itemize}
\item $\partial^\alpha h_\beta(y)=\sqrt{\alpha!/\gamma!}h_\gamma(y)$, where $\gamma=\beta-\alpha$,
\item $h_\beta(\sqrt{1-\lambda}x+\sqrt\lambda y)=(1-\lambda)^{|\beta|/2}\sum_{\alpha\leq\beta}\frac{\partial^\alpha h_\beta(x)}{\sqrt{\alpha!}}\left(\frac\lambda{1-\lambda}\right)^{|\alpha|/2}h_\alpha(y)$,
\item $\partial^\alpha h_\beta(\sqrt{1-\lambda}x+\sqrt\lambda y)=(1-\lambda)^{|\beta-\alpha|/2}\sum_{\gamma\leq\beta-\alpha}\frac{\partial^{\alpha+\gamma}h_\beta(x)}{\sqrt{\gamma!}}\left(\frac\lambda{1-\lambda}\right)^{|\gamma|/2}h_\gamma(y)$.
\end{itemize}
We conclude with a Taylor-like expansion in the Hermite basis that we use repeatedly.
\paragraph{Lemma 16.} Let $f(y)=\sum_\alpha\widehat f(\alpha)h_\alpha(y)$. Then
\[f(\sqrt{1-\lambda}x+\sqrt\lambda y)=\sum_\alpha\frac{\partial^\alpha g(x)}{\sqrt{\alpha!}}\left(\frac\lambda{1-\lambda}\right)^{|\alpha|/2}h_\alpha(y),\]
where $g(x):=U_{\sqrt{1-\lambda}}f(x)=\sum_\alpha\widehat f(\alpha)(1-\lambda)^{|\alpha|/2}h_\alpha(x)$.
\paragraph{Proof.} We express
\begin{align*}
f(\sqrt{1-\lambda}x+\sqrt\lambda y)&=\sum_\alpha\widehat f(\alpha)h_\alpha(\sqrt{1-\lambda}x+\sqrt\lambda y)\\
&=\sum_\alpha\frac{h_\alpha(y)}{\sqrt{\alpha!}}\left(\frac\lambda{1-\lambda}\right)^{|\alpha|/2}\sum_{\beta\geq\alpha}\widehat f(\beta)(1-\lambda)^{|\beta|/2}\partial^\alpha h_\beta(x)\\
&=\sum_\alpha\frac{h_\alpha(y)}{\sqrt{\alpha!}}\left(\frac\lambda{1-\lambda}\right)^{|\alpha|/2}\partial^\alpha g(x).
\end{align*}\hfill$\triangleleft$
Lastly, we will also need an extension of this theorem which expresses $\partial^\alpha f$, at the point $\sqrt{1-\lambda}x+\sqrt\lambda y$, as a polynomial in $y$ in the Hermite basis.
\paragraph{Theorem 17.} Let $f(y)=\sum_\alpha\widehat f(\alpha)h_\alpha(y)$. Then
\[\partial^\alpha f(\sqrt{1-\lambda}x+\sqrt\lambda y)=(1-\lambda)^{-|\alpha|/2}\sum_{\beta\geq\alpha}\partial^\beta g(x)\sqrt{\frac{\alpha!}{\beta!}}\left(\frac\lambda{1-\lambda}\right)^{|\beta-\alpha|/2}h_{\beta-\alpha}(y),\]
where $g(x):=U_{\sqrt{1-\lambda}}f(x)$.
\paragraph{Proof.} We express
\begin{align*}
\partial^\alpha f(\sqrt{1-\lambda}x+\sqrt\lambda y)&=\sum_\beta\widehat f(\beta)\partial^\alpha h_\beta(\sqrt{1-\lambda}x+\sqrt\lambda y)\\
&=\sum_\gamma\frac{h_\gamma(y)}{\sqrt{\gamma!}}\left(\frac\lambda{1-\lambda}\right)^{|\gamma|/2}\sum_{\beta\geq\gamma+\alpha}\widehat f(\beta)(1-\lambda)^{|\beta-\alpha|/2}\partial^{\alpha+\gamma}h_\beta(x)\\
&=(1-\lambda)^{-|\alpha|/2}\sum_\gamma\frac{h_\gamma(y)}{\sqrt{\gamma!}}\left(\frac\lambda{1-\lambda}\right)^{|\gamma|/2}\partial^{\alpha+\gamma}g(x).
\end{align*}\hfill$\triangleleft$

\section*{4.1 Proof of Lemma 6}
Our main tool will be Kane's relative-anticoncentration lemma for gaussian polynomials.
\paragraph{Lemma 18 ([9]).} For a degree $d$ polynomial $p$, and independent standard gaussian vectors $x,y\in\mathbb R^n$,
\[\mathbb P(|p(x)|\leq\varepsilon|\langle y,\nabla p(x)\rangle|)\leq O(\varepsilon d^2).\]
In fact, we will actually work with the following corollary which is essentially the first of the $d$ inequalities in Lemma 6.
\paragraph{Corollary 19.} For a degree $d$ polynomial $p$, and independent standard gaussian vector $x\in\mathbb R^n$,
\[\mathbb P(|p(x)|\leq\varepsilon\|\nabla p(x)\|)\leq O(\varepsilon d^2).\]
\paragraph{Proof.} We note that for any fixed $x$, $\langle y,\nabla p(x)\rangle$ is identical in distribution to $Z\|\nabla p(x)\|$, where $Z\sim N(0,1)$ is a standard gaussian. So, we express
\begin{align*}
\mathbb P(|p(x)|\leq\varepsilon|\langle y,\nabla p(x)\rangle|)&=\mathbb P(|p(x)|\leq\varepsilon|Z|\|\nabla p(x)\|)\\
&\geq\mathbb P(|p(x)|\leq\varepsilon\|\nabla p(x)\|)\cdot\mathbb P(|Z|\geq1).
\end{align*}
Since $\mathbb P(|Z|\geq1)\geq\Omega(1)$, we conclude that
\[\mathbb P(|p(x)|\leq\varepsilon\|\nabla p(x)\|)\leq O(\varepsilon d^2).\]\hfill$\triangleleft$
The heart of the proof of Lemma 6 is a vector-valued variant of the above corollary:
\paragraph{Lemma 20.} Let $\vec f(x):=(f_1(x),f_2(x),\ldots,f_m(x))$ be a collection of $m$ degree-at-most $d$ polynomials $f_j(x)$. If $x\in\mathbb R^n$ is a standard gaussian vector, then
\[\mathbb P\left(\|\vec f(x)\|^2\leq\varepsilon^2\sum_{j=1}^m\|\nabla f_j(x)\|^2\right)\leq O(\varepsilon d^2).\]
\paragraph{Proof of Lemma 6.} We simply apply the above lemma $d$ times and take a union bound. For $1\leq k\leq d$, let $\vec f_k(x):=((\partial^\alpha f(x):|\alpha|=k))$. Note that $\|\vec f_k(x)\|^2=\|\nabla^kf(x)\|^2$. Further, note that
\[\sum_{\alpha:|\alpha|=k}\|\nabla(\partial^\alpha f(x))\|^2\geq\|\nabla^{k+1}f(x)\|^2,\]
where the inequality follows as each $(k+1)$'th order derivative would be counted at least once in the expression on the left hand side. Therefore, by the above lemma, for $x\sim N(0,1)^n$, we have
\[\mathbb P(\|\nabla^kf(x)\|^2\leq\varepsilon^2\|\nabla^{k+1}f(x)\|^2)\leq O(\varepsilon d^2).\]
Setting $\varepsilon=\delta/d^3$, and taking a union bound over all $k$, we get that for a constant $C>0$,
\[\mathbb P(\forall k,\|\nabla^kf(x)\|^2>C(\delta^2/d^6)\|\nabla^{k+1}f(x)\|^2)\geq1-\delta.\]
This proves Lemma 6.\hfill$\triangleleft$
\paragraph{Proof of Lemma 20.} Consider the auxiliary polynomial
\[h(x,y):=\sum_{j=1}^mf_j(x)y_j.\]
As a function of both $x$ and $y$, we have
\[\nabla h(x,y)=\vec f(x)\circ M_xy,\]
where $M_x$ is the matrix with columns $\nabla f_j(x)$ (that is, $M_x$ has $(i,j)$-th entry $\frac\partial{\partial x_i}f_j(x)$). So, applying Corollary 19 to this auxiliary polynomial gives the probability bound
\begin{align*}
q&:=\mathbb P(h(x,y)^2\leq\varepsilon^2\|\nabla g(x,y)\|^2)\\
&=\mathbb P(\langle y,\vec f(x)\rangle^2\leq\varepsilon^2(\|\vec f(x)\|^2+\|M_xy\|^2))\\
&\leq O(\varepsilon d^2).
\end{align*}
Now, for some constant $C\geq2$ to be specified later, let $E$ denote the event that
\[(C^2-1)\|\vec f(x)\|^2\leq\frac{\varepsilon^2}2\|M_x\|_F^2,\]
where $\|M_x\|_F$ is the Frobenius norm of $M_x$. We note that we can lower-bound the probability $q$ by
\[q\geq\mathbb P(E)\cdot\mathbb P(\langle y,\vec f(x)\rangle\leq C\|\vec f(x)\|\text{ and }\|M_xy\|^2\geq\tfrac12\|M_x\|_F^2\mid E).\]
We claim that for large enough choice of constant $C$, this conditional probability can be lower-bounded by $\Omega(1)$. Indeed, we can argue for any fixed $x$:
\begin{itemize}
\item $\mathbb P(\langle y,\vec f(x)\rangle\geq C\|\vec f(x)\|)\leq1/C^2$.
\item $\mathbb P(\|M_xy\|^2\geq\tfrac12\|M_x\|_F^2)\geq\Omega(1)$.
\end{itemize}
The first item is just a Chebyshev inequality; the second item can be derived e.g. from the basic anticoncentration bound one obtains for degree-$2$ polynomials from the Paley-Zygmund bound together with hypercontractivity (since, for any fixed matrix $M$, the quadratic form $g(y):=\|My\|^2$ has second-moment $\mathbb Eg(y)^2\geq(\mathbb Eg(y))^2=\|M\|_F^2$).
Thus, by choosing $C$ large enough, we can lower-bound this conditional probability by $\Omega(1)-1/C^2\geq\Omega(1)$.
We conclude that $\mathbb P(E)\leq O(q)=O(\varepsilon d^2)$. This gives the desired conclusion
\[\mathbb P(\|\vec f(x)\|\leq\Omega(\varepsilon)\|M_x\|_F)\leq O(\varepsilon d^2).\]\hfill$\triangleleft$
\section*{5 Pseudorandom Generator for PTFs}
The following theorem gives quantitative bounds on the error of our main generator:
\paragraph{Theorem 21.} Fix some parameters $\varepsilon>0$ and $R\in\mathbb N$. Let $z$ be a standard gaussian, and let $Z=\frac1{\sqrt L}\sum_{i=1}^LY_i$, where each $Y_i$ is $dR$-moment-matching. Then for some sufficiently large absolute constant $c$ and any polynomial $p$ of degree $d$,
\[\mathbb E_Z\operatorname{sign}(p(Z))\geq\mathbb E_{z\sim N(0,1)^n}\operatorname{sign}(p(z))-O(\varepsilon d^3)-L\cdot2^{-\Omega(R)},\]
as long as $L$ is at least $Rd^c/\varepsilon^2$.
Combining the above with Theorem 12 immediately implies our main result Theorem 3.
\paragraph{Proof of Theorem 3.} Given a target error $\varepsilon'$, set $\varepsilon=\varepsilon'/Cd^3$, and $R=C\log(d/\varepsilon)$ for a sufficiently big constant so that the error in the above lemma is at most $\varepsilon'/2$ for $L=Rd^c/\varepsilon^2=O(d^c\log(d/\varepsilon)/\varepsilon^2)$. While the above theorem only gives a lower bound, we can get an upper bound by applying the result to $-p$. Now, by applying Theorem 12 there exists an efficient PRG that fools degree $d$ PTFs with error at most $\varepsilon'$ and seed length $O(d^{O(1)}\log(nd/\varepsilon')\log(d/\varepsilon')/(\varepsilon')^2)$ which can be simplified to the bound in the theorem.\hfill$\triangleleft$
We now prove the above theorem by the lower-sandwiching argument outlined in Section 2.1. Fix a polynomial $p(x)$ of degree $d$. We remind the reader of our convention $\operatorname{sign}(t):=\boldsymbol1(t\geq0)$.
We define the mollifier function
\[g(x):=\prod_{k=0}^{d-1}\rho\left(\log\left(\frac1{16\varepsilon^2}\frac{\|\nabla^kp(x)\|^2}{\|\nabla^{k+1}p(x)\|^2}\right)\right),\]
where $\rho:\mathbb R\to[0,1]$ is some smooth univariate function with $\rho(t)=0$ for $t\leq0$, $\rho(t)=1$ for $t\geq1$, and $\|\frac{\partial^k\rho}{\partial t^k}\|_\infty\leq k^{O(k)}$ for all $k$.\footnote{For example, it suffices to let $\rho(t)$ be the standard mollifier $\rho(t):=0$ for $t\leq0$, $\rho(t):=1$ for $t\geq1$, and $\rho(t):=e\cdot\exp(1/((t-1)^2-1))$ for $t\in(0,1)$.}
\paragraph{Proof of Theorem 21.} For every point $x\in\mathbb R^n$ we have
\[\operatorname{sign}(p(x))\geq\operatorname{sign}(p(x))g(x).\]
Furthermore, under the truly-random gaussian inputs $z\sim N(0,1)^n$ we have
\[\mathbb E_z\operatorname{sign}(p(z))g(z)\geq\mathbb E_z\operatorname{sign}(p(z))-\mathbb E_z|g(z)-1|\geq\mathbb E_z\operatorname{sign}(p(z))-O(\varepsilon d^3),\]
where the final inequality here follows from Lemma 6. Combining these, we get that
\begin{multline*}
\mathbb E_Z\operatorname{sign}(p(Z))\geq\mathbb E_{z\sim N(0,1)^n}\operatorname{sign}(p(z))-O(\varepsilon d^3)\\
-\left|\mathbb E_Z\operatorname{sign}(p(Z))g(Z)-\mathbb E_{z\sim N(0,1)^n}\operatorname{sign}(p(z))g(z)\right|.
\end{multline*}
Thus, it suffices to bound $|\mathbb E_Z\operatorname{sign}(p(Z))g(Z)-\mathbb E_{z\sim N(0,1)^n}\operatorname{sign}(p(z))g(z)|$, which we do by a hybrid argument. We first represent $z$ as $z:=\frac1{\sqrt L}\sum_{i=1}^Ly_i$ where each $y_i$ is an independent standard gaussian. We can replace each $Y_i$ with $y_i$ and get
\[|\mathbb E_Z\operatorname{sign}(p(Z))g(Z)-\mathbb E_y\operatorname{sign}(p(y))g(y)|\leq2^{-\Omega(R)}L,\]
as a consequence of the following lemma (restatement of Lemma 8) that we prove in the next section. Theorem 21 now follows.\hfill$\triangleleft$
\paragraph{Lemma 22 (Main hybrid-step).} There exists a constant $c$ such that the following holds for $\lambda\leq\varepsilon^2/Rd^c$. For any fixed vector $x\in\mathbb R^n$, $Y$ a $dR$-moment-matching gaussian vector, and $y\sim N(0,1)^n$,
\[|\mathbb E_Y\operatorname{sign}(p(x+\sqrt\lambda Y))g(x+\sqrt\lambda Y)-\mathbb E_y\operatorname{sign}(p(x+\sqrt\lambda y))g(x+\sqrt\lambda y)|\leq\gamma=2^{-\Omega(R)}.\]
\subsection*{5.1 Analysis of the main hybrid-step}
The proof of Lemma 22 is by a case-analysis as outlined in the introduction. Consider the setting as in the lemma and define
\[\varphi(z):=U_{\sqrt{1-\lambda}}p\left(\frac z{\sqrt{1-\lambda}}\right).\]
The core argument will be a case-analysis on the derivatives of $\varphi$ at the fixed point $x$ and whether these are slow-growing. Note that if $p$ were multi-linear, then we would simply have $\varphi\equiv p$. The starting point is the following re-scaling of Lemma 16:
\[p(x+\sqrt\lambda y)=\sum_{|\alpha|\leq d}\frac{\partial^\alpha\varphi(x)}{\sqrt{\alpha!}}\lambda^{|\alpha|/2}h_\alpha(y).\tag{4}\]
Further, by a re-scaling of Theorem 17, we get the following identity which gives a nice nearly self-referential expression relating the derivatives of $p$ to those of $\varphi$:
\[\partial^\alpha p(x+\sqrt\lambda y)=\sum_{\beta\geq\alpha}\sqrt{\frac{\alpha!}{\beta!}}\partial^\beta\varphi(x)\lambda^{|\beta-\alpha|/2}h_{\beta-\alpha}(y).\tag{5}\]
Now, note that for a truly random gaussian $y$ we have $\partial^\alpha\varphi(x)=\mathbb E_y\partial^\alpha p(x+\sqrt\lambda y)$. Thus, it is reasonable to expect that for typical points $x$ and small enough $\lambda$, $\partial^\alpha p(x+\sqrt\lambda y)$ will be strongly concentrated around $\partial^\alpha\varphi(x)$. The following lemma gives quantitative bounds on how much the derivatives $\partial^\alpha p(x+\sqrt\lambda y)$ deviate from their expectations $\partial^\alpha\varphi(x)$ for a random $y\sim N(0,1)^n$. As we will need such bounds even for $k$-moment-matching $Y$, we state the deviation bound in terms of moments:
\paragraph{Lemma 23.} Suppose $f$ is a degree-$d$ polynomial, and let $\varphi(z)=U_{\sqrt{1-\lambda}}f(z/\sqrt{1-\lambda})$. Consider the polynomial
\[D(y):=\|\nabla^kf(x+\sqrt\lambda y)-\nabla^k\varphi(x)\|^2,\]
which measures the euclidean distance between the $k$-th order derivatives $\nabla^kf(x+\sqrt\lambda y)$ and their expectations $\nabla^k\varphi(x)$. For $y\sim N(0,1)^n$, we have the moment bound
\[\|D(y)\|_{q/2}\leq\sum_{t=k+1}^d(\lambda dq)^{t-k}\|\nabla^t\varphi(x)\|^2.\]
That is,
\[\left(\mathbb E_{y\sim N(0,1)^n}\|\nabla^kf(x+\sqrt\lambda y)-\nabla^k\varphi(x)\|^q\right)^{1/q}\leq\sqrt{\sum_{t=k+1}^d(\lambda dq)^{t-k}\|\nabla^t\varphi(x)\|^2}.\]
\paragraph{Proof.} We express
\begin{align*}
D(y)&=\sum_\alpha(\partial^\alpha f(x+\sqrt\lambda y)-\partial^\alpha\varphi(x))^2\\
&=\sum_\alpha\left(\sum_{\beta>\alpha}\sqrt{\frac{\alpha!}{\beta!}}\partial^\beta\varphi(x)\lambda^{|\beta-\alpha|/2}h_{\beta-\alpha}(y)\right)^2.
\end{align*}
First, by triangle-inequality, we get
\begin{align*}
\|D(y)\|_{q/2}&\leq\sum_\alpha\left\|\left(\sum_{\beta>\alpha}\sqrt{\frac{\alpha!}{\beta!}}\partial^\beta\varphi(x)\lambda^{|\beta-\alpha|/2}h_{\beta-\alpha}(y)\right)^2\right\|_{q/2}\\
&=\sum_\alpha\left\|\sum_{\beta>\alpha}\sqrt{\frac{\alpha!}{\beta!}}\partial^\beta\varphi(x)\lambda^{|\beta-\alpha|/2}h_{\beta-\alpha}(y)\right\|_q^2.
\end{align*}
Applying hypercontractivity, we now get
\begin{align*}
\|D(y)\|_{q/2}&\leq\sum_\alpha\left\|U_{\sqrt q}\sum_{\beta>\alpha}\sqrt{\frac{\alpha!}{\beta!}}\partial^\beta\varphi(x)\lambda^{|\beta-\alpha|/2}h_{\beta-\alpha}(y)\right\|_2^2\\
&=\sum_\alpha\sum_{\beta>\alpha}\frac{\alpha!}{\beta!}\partial^\beta\varphi(x)^2\lambda^{|\beta-\alpha|}q^{|\beta-\alpha|}\\
&\leq\sum_\alpha\sum_{\beta>\alpha}\partial^\beta\varphi(x)^2\lambda^{|\beta-\alpha|}q^{|\beta-\alpha|}\\
&=\sum_{t=k+1}^d\binom t{t-k}(\lambda q)^{t-k}\|\nabla^t\varphi(x)\|^2\\
&\leq\sum_{t=k+1}^d(\lambda dq)^{t-k}\|\nabla^t\varphi(x)\|^2.
\end{align*}\hfill$\triangleleft$
We are now ready to prove Lemma 22.
\paragraph{Proof of Lemma 22.} We study two cases:
\begin{enumerate}
\item $x$ is poorly-behaved for $\varphi$. In this case, we will show that $g(x+\sqrt\lambda Y)=0$ with probability at least $1-2^{-\Omega(R)}$.
\item $x$ is well-behaved for $\varphi$: In this case, we will exploit the fact that $\operatorname{sign}(p(x+\sqrt\lambda Y))$ will equal $\operatorname{sign}(\varphi(x))$ with probability $1-2^{-\Omega(R)}$. We then have to show that $Y$ fools the mollifier $g$ which is a bit technically involved (hence we deal with this case second unlike in Section 2.1).
\end{enumerate}
We begin with the first case.
\paragraph{Case 1: $x$ is poorly-behaved for $\varphi$.} Consider the case where the inequality $\|\nabla^k\varphi(x)\|\geq\varepsilon\|\nabla^{k+1}\varphi(x)\|$ is violated for some $k$, and indeed let $k_0$ be the largest $k$ such that this inequality is violated. We will argue that with probability at least $1-2^{-\Omega(R)}$, over random choice of $Y$, that
\[\|\nabla^{k_0}p(x+\sqrt\lambda Y)\|\leq4\varepsilon\|\nabla^{k_0+1}p(x+\sqrt\lambda Y)\|,\]
in which case $g(x+\sqrt\lambda Y)=0$.
More specifically, we will show that it is highly likely that both
\begin{itemize}
\item $\|\nabla^{k_0}p(x+\sqrt\lambda Y)\|\leq2\varepsilon\|\nabla^{k_0+1}\varphi(x)\|$, and
\item $\|\nabla^{k_0+1}p(x+\sqrt\lambda Y)\|\geq\frac12\|\nabla^{k_0+1}\varphi(x)\|$.
\end{itemize}
For this, we will use Equation (5) and Lemma 23. Supposing $k_0$ is the largest $k$ such that $\|\nabla^k\varphi(x)\|<\varepsilon\|\nabla^{k+1}\varphi(x)\|$, we have
\begin{itemize}
\item $\|\nabla^{k_0}\varphi(x)\|\leq\varepsilon\|\nabla^{k_0+1}\varphi(x)\|$ and
\item $\|\nabla^{k_0+1}\varphi(x)\|\geq\varepsilon^t\|\nabla^{k_0+1+t}\varphi(x)\|$ for all $t\geq0$.
\end{itemize}
Lemma 23 therefore gives the bounds
\[\left(\mathbb E_Y\|\nabla^{k_0}p(x+\sqrt\lambda Y)-\nabla^{k_0}\varphi(x)\|^R\right)^{1/R}\leq\varepsilon\|\nabla^{k_0+1}\varphi(x)\|\sqrt{\sum_{t\geq1}(\lambda dR/\varepsilon^2)^t}\]
and
\[\left(\mathbb E_Y\|\nabla^{k_0+1}p(x+\sqrt\lambda Y)-\nabla^{k_0+1}\varphi(x)\|^R\right)^{1/R}\leq\|\nabla^{k_0+1}\varphi(x)\|\sqrt{\sum_{t\geq1}(\lambda dR/\varepsilon^2)^t}.\]
So, as long as $\lambda dR/\varepsilon^2$ is at most a sufficiently small constant, we conclude that the following bounds hold with probability at least $1-2^{-R}$:
\begin{align*}
\|\nabla^{k_0}p(x+\sqrt\lambda Y)\|&\leq\|\nabla^{k_0}\varphi(x)\|+\|\nabla^{k_0}p(x+\sqrt\lambda Y)-\nabla^{k_0}\varphi(x)\|\\
&\leq2\varepsilon\|\nabla^{k_0+1}\varphi(x)\|,\text{ and}\\
\|\nabla^{k_0+1}p(x+\sqrt\lambda Y)\|&\geq\|\nabla^{k_0+1}\varphi(x)\|-\|\nabla^{k_0+1}p(x+\sqrt\lambda Y)-\nabla^{k_0+1}\varphi(x)\|\\
&\geq\tfrac12\|\nabla^{k_0+1}\varphi(x)\|.
\end{align*}
In the case that these bounds hold, we get
\[\|\nabla^{k_0}p(x+\sqrt\lambda Y)\|\leq4\varepsilon\|\nabla^{k_0+1}p(x+\sqrt\lambda Y)\|,\]
and so $g(x+\sqrt\lambda Y)=0$. As this holds with probability at least $1-2^{-\Omega(R)}$ for both $y\sim N(0,1)^n$ as well as $Y$, the conclusion of Lemma 22 follows. This finishes the proof of Case 1.\hfill$\lrcorner$
\paragraph{Case 2: $x$ is well-behaved for $\varphi$.} We now consider the complimentary case where
\[\|\nabla^k\varphi(x)\|\geq\varepsilon\|\nabla^{k+1}\varphi(x)\|\]
for all $k=0,1,\ldots,d-1$. Consider the normalized polynomial
\[f(y):=\frac{p(x+\sqrt\lambda y)}{\varphi(x)}=1+\frac1{\varphi(x)}\sum_{\alpha\neq0}\partial^\alpha\varphi(x)\lambda^{|\alpha|/2}h_\alpha(y).\]
Using hypercontractivity, we bound the $R$-th moment of $f(y)-1$ by its $\sqrt R$-hypervariance:
\[\|f(y)-1\|_R\leq\|U_{\sqrt R}(f(y)-1)\|_2\leq\sqrt{\sum_{k\geq1}(\lambda R/\varepsilon^2)^k}\leq1/2.\]
So, by a Markov argument, we have
\[\mathbb P(\operatorname{sign}(p(x+\sqrt\lambda Y))\neq\operatorname{sign}(\varphi(x)))\leq2^{-R},\]
and this holds whenever $Y$ is $k$-moment-matching for $k\geq dR$. So, $\operatorname{sign}(p(x+\sqrt\lambda Y))$ is nearly a constant for random $Y$; it remains to show that $Y$ fools $g(x+\sqrt\lambda Y)$. We do this by (essentially) truncating the Taylor-series of $g$ about $x$ so that we are left with a degree $dR$ polynomial, which is fooled by $Y$. The truncation-error will be small because our assumption, $\|\nabla^k\varphi(x)\|\geq\varepsilon\|\nabla^{k+1}\varphi(x)\|$ for all $k$, gives us good control on the $R$-th order moments of the deviations $\|\nabla^k\varphi(x)\|-\|\nabla^kp(x+\sqrt\lambda Y)\|$. The exact calculations are somewhat cumbersome and are given below. We will show that $Y$ fools the mollifier function
\[g(x+\sqrt\lambda y)=\prod_{k=0}^{d-1}\rho\left(\log\left(\frac1{16\varepsilon^2}\frac{\|\nabla^kp(x+\sqrt\lambda y)\|^2}{\|\nabla^{k+1}p(x+\sqrt\lambda y)\|^2}\right)\right).\]
To simplify notation we define the shifted function $\sigma(t):=\rho(t-\log(16\varepsilon^2))$, and express
\[g(x+\sqrt\lambda y)=\prod_{k=0}^{d-1}\sigma(\log\|\nabla^kp(x+\sqrt\lambda y)\|^2-\log\|\nabla^{k+1}p(x+\sqrt\lambda y)\|^2).\]
It will be convenient to think of $g$ (redundantly) as function of $2d$ auxiliary variables $s_1\ldots s_d,t_1,\ldots,t_d$, which we will eventually fix to
\begin{itemize}
\item $s_i:=\|\nabla^{i-1}p(x+\sqrt\lambda y)\|^2$
\item $t_i:=\|\nabla^ip(x+\sqrt\lambda y)\|^2$,
\end{itemize}
so we write $g(s,t):=\prod_{i=1}^d\sigma(\log(s_i)-\log(t_i))$.
We Taylor-expand $g(s,t)$ around the points
\begin{itemize}
\item $a_i:=\|\nabla^{i-1}\varphi(x)\|^2$
\item $b_i:=\|\nabla^i\varphi(x)\|^2$,
\end{itemize}
which gives $g(s,t)=\ell(s,t)+h(s,t)$, with low-degree part
\[\ell(s,t):=\sum_{\substack{\alpha,\beta\in\mathbb N^d\\|\alpha|+|\beta|<R}}\frac{\partial_s^\alpha\partial_t^\beta g(a,b)}{\alpha!\beta!}(s-a)^\alpha(t-b)^\beta\]
and remainder
\[|h(s,t)|\leq\sum_{\substack{\alpha,\beta\in\mathbb N^d\\|\alpha|+|\beta|=R}}\frac{|\partial_s^\alpha\partial_t^\beta g(s^*,t^*)|}{\alpha!\beta!}|s-a|^\alpha|t-b|^\beta,\]
where ``$|\partial_s^\alpha\partial_t^\beta g(s^*,t^*)|$'' is notation for the maximum magnitude of $\partial_s^\alpha\partial_t^\beta g$ on any point on the line segment from $(a,b)$ to $(s,t)$. We need the following fact to bound the size of the derivatives of $g$,
\paragraph{Claim 24.} Suppose $\sigma$ is a smooth univariate function with uniform derivative bounds $\|\sigma^{(n)}\|_\infty\leq n^{O(n)}$. The bivariate function $r(u,v):=\sigma(\log(u)-\log(v))$ has derivatives bounded in size by
\[\left|\frac{\partial^n}{\partial u^n}\frac{\partial^m}{\partial v^m}r(u,v)\right|\leq\frac{n^{O(n)}m^{O(m)}}{|u|^n|v|^m}.\]
This claim follows easily from the generalized chain rule (Faà di Bruno's formula). As a result, we get the derivative bounds
\[|\partial_s^\alpha\partial_t^\beta g(s,t)|\leq\frac{|\alpha|^{O(|\alpha|)}}{|s^\alpha|}\frac{|\beta|^{O(|\beta|)}}{|t^\beta|}.\]
Using this, we bound the remainder
\[|h(s,t)|\leq\sum_{\substack{\alpha,\beta\in\mathbb N^d\\|\alpha|+|\beta|=R}}d^{O(R)}\prod_{i=1}^d\left(\frac{|1-s_i/a_i|}{1-|1-s_i/a_i|}\right)^{\alpha_i}\left(\frac{|1-t_i/b_i|}{1-|1-t_i/b_i|}\right)^{\beta_i}.\]
Now, consider the event $E$ (which depends on $y$) that
\[(1-\delta)\|\nabla^i\varphi(x)\|^2\leq\|\nabla^ip(x+\sqrt\lambda y)\|^2\leq(1+\delta)\|\nabla^i\varphi(x)\|^2\]
holds for all $i$, where $\delta\leq1/2$ is a parameter we will set shortly. In the case that this indeed holds, we get $|h(s,t)|\leq d^{O(R)}O(\delta)^R$.
We set $\delta$ just small enough to ensure $|h(s,t)|\leq2^{-R}$.
Now, we express $g$ (which we now think of as a function of the underlying variable $y$) as
\begin{align*}
g&=g\cdot\boldsymbol1_E+g\cdot\boldsymbol1_{\overline E}\\
&=\ell\cdot\boldsymbol1_E+h\cdot\boldsymbol1_E+g\cdot\boldsymbol1_{\overline E}\\
&=\ell-\ell\cdot\boldsymbol1_{\overline E}+h\cdot\boldsymbol1_E+g\cdot\boldsymbol1_{\overline E},
\end{align*}
and we obtain the pointwise bound
\[|g-\ell|\leq2^{-R}+\boldsymbol1_{\overline E}+|\ell|\cdot\boldsymbol1_{\overline E}.\]
On average over $Y$, we get truncation error
\begin{align*}
\mathbb E_Y|g(x+\sqrt\lambda Y)-\ell(Y)|
&\leq2^{-R}+\mathbb E_Y\boldsymbol1_{\overline E}(Y)+\sqrt{\mathbb E_Y\ell^2(Y)}\sqrt{\mathbb E_Y\boldsymbol1_{\overline E}(Y)}\\
&\leq2^{-R}+O(d/\delta)^R\cdot(\lambda dR/\varepsilon^2)^{-\Omega(R)}\\
&\leq2^{-R}+d^{O(1)}\cdot(\lambda dR/\varepsilon^2)^{-\Omega(R)}
\end{align*}
where the second inequality here follows from the moment bounds in Lemma 23. As required by the conditions of Lemma 22, we insist that $\lambda$ is small enough that this error is at most $2^{-\Omega(R)}$. Since this bound holds also for truly-random standard gaussian $y$, and $\mathbb E_Y\ell(Y)=\mathbb E_y\ell(y)$, we obtain the desired bound
\[|\mathbb E_Yg(x+\sqrt\lambda Y)-\mathbb E_yg(x+\sqrt\lambda y)|\leq2^{-\Omega(R)}.\]
This finishes the proof in Case 2 and hence of Lemma 22.\hfill$\triangleleft$
